\section{Checklist de conception}
\label{sec:checklist_conception}

\begin{UPSTIidee}{À décider avant de commencer le TP}
    Avant d'ouvrir Codesys, formalisez par écrit les points suivants pour votre système. Cette fiche peut être complétée directement et conservée comme mémo de conception.
\end{UPSTIidee}

\begin{enumerate}
    \item \textbf{Rôles} : quel équipement est Scanner ? Quel(s) équipement(s) sont Adapter(s) ? (\ref{sec:architecture_reseau})
    \item \textbf{Adressage} : quelle adresse IP pour chaque équipement ? Sont-elles sur le même sous-réseau ?
    \item \textbf{Données échangées} : lister précisément, pour chaque Adapter, les données à transmettre du Scanner vers l'Adapter (commandes) et de l'Adapter vers le Scanner (états/mesures).
    \item \textbf{Dimensionnement des Assembly} : taille en octets de l'Input Assembly, de l'Output Assembly, et de la Configuration Assembly le cas échéant. (\ref{sec:objets_assembly})
    \item \textbf{RPI} : quelle période de rafraîchissement cyclique choisir pour chaque connexion, en cohérence avec la dynamique du procédé et la charge réseau tolérable ? (\ref{sec:messagerie})
    \item \textbf{Messagerie explicite} : quelles informations ne nécessitent qu'une lecture/écriture ponctuelle (identification, diagnostic, paramétrage) plutôt qu'un échange cyclique ?
    \item \textbf{Gestion des erreurs de connexion} : comment le programme détecte-t-il une perte de connexion avec un Adapter (bit de statut du scanner, timeout) ? Quel comportement de repli adopter (arrêt sécurisé, dernière valeur connue, signalisation) ?
    \item \textbf{Découpage objet} : quelle classe (ou quel ensemble de classes) va encapsuler la communication avec chaque équipement ? Quelles propriétés et méthodes seront exposées au reste de l'application ? (\ref{sec:mise_en_oeuvre_codesys})
    \item \textbf{Documentation constructeur} : le fichier EDS de l'équipement est-il disponible ? Quels objets/Assembly spécifiques au métier de l'équipement propose-t-il, au-delà des objets standards ?
\end{enumerate}

\begin{UPSTIinfor}{Erreurs fréquentes à anticiper}
    \begin{itemize}
        \item Confondre le sens Input/Output entre le point de vue de l'Adapter (documentation constructeur) et celui du Scanner (votre programme).
        \item Ne pas gérer explicitement la perte de connexion : un Adapter déconnecté peut laisser des données figées sur leur dernière valeur si aucune détection n'est mise en place.
        \item Sous-dimensionner ou surdimensionner un Assembly par rapport aux besoins réels, rendant le mapping confus ou gaspillant de la bande passante.
        \item Placer la logique métier directement sur les variables brutes du mapping, au lieu de l'encapsuler dans une classe dédiée -- ce qui rend le code difficile à faire évoluer si l'équipement change.
    \end{itemize}
\end{UPSTIinfor}
